% TEMPLATE-MILC-v3.tex - plantilla maestra UNICA de las guias MILC v3.
%
% La usan las tres familias de guias, cada una con su driver:
%   · Grados 6-11  -> scripts/build-guias-g11.py
%   · Bebras       -> scripts/build-guias-bebras.py
%   · Semillero    -> scripts/build-guias-semillero.py
%
% Antes habia tres copias casi identicas (~400 lineas iguales, 6 distintas):
% cada mejora habia que aplicarla tres veces, y en la practica no se hacia
% (el motor de recursos vivia solo en dos de ellas). Lo que varia entre
% familias esta parametrizado en 6 placeholders que cada driver llena
% (se nombran aqui SIN su sintaxis real: el builder reemplaza por texto plano,
% tambien dentro de los comentarios, y un valor multilinea rompe el preambulo):
%
%   HEADER_IZQ        encabezado izquierdo de pagina
%   HEADER_DER        encabezado derecho de pagina
%   PORTADA_ETIQUETA  rotulo sobre el numero grande (GRADO/MOMENTO/MODULO)
%   PORTADA_SERIE     linea de serie de la portada
%   PORTADA_ID        identificador de la guia en la portada
%   PORTADA_CIRCULO   el circulito de grado (°), o vacio si no aplica
%
% El resto de claves las documenta content/guias/_SCHEMA.md. El script valida
% que no queden placeholders sin reemplazar antes de escribir el archivo final.

\documentclass[11pt,letterpaper]{article}
\usepackage[margin=1.75cm]{geometry}
\usepackage[table]{xcolor}
\usepackage{fontspec}
\usepackage[spanish]{babel}
\usepackage{tikz}
% Librerías TikZ para los diagramas de las guías (bloque `recursos:`):
%  · babel  → babel en español vuelve activos caracteres como «>», y TikZ los usa
%             en sus opciones (>=stealth, ->). Sin esta librería, cualquier
%             diagrama con flechas revienta con «\language@active@arg> has an
%             extra }». Debe ir ANTES de que se dibuje nada.
%  · positioning        → sintaxis «below left=2pt and 3pt».
%  · decorations.pathreplacing → llaves ({brace}) para señalar rangos.
\usetikzlibrary{babel,positioning,decorations.pathreplacing}
\usepackage{graphicx}
% Circuitos: el trabajo de electrónica se hace hoy con micro:bit y sus sensores
% integrados (MakeCode, sin cableado), así que no se necesitan esquemáticos.
% Si algún día entran montajes con protoboard, basta descomentar esta línea:
% los diagramas .tex/.tikz declarados en `recursos:` ya se incrustan solos.
% \usepackage{circuitikz}
\usepackage[most]{tcolorbox}
\usepackage{tabularx}
\usepackage{array}
\usepackage{fancyhdr}
\usepackage{titlesec}
\usepackage{ragged2e}
\usepackage{enumitem}
\usepackage{microtype}

% Helvetica Neue no trae la flecha U+2192, asi que cada «→» escrito literalmente
% en el YAML se compilaba a NADA, en silencio (462 flechas en 61 guias: rutas del
% tipo «paleta → tipografia → lockup» salian rotas). Mapeamos el caracter a su
% version matematica, que si tiene glifo. Asi el autor puede seguir escribiendo
% «→» comodamente en el YAML.
\usepackage{newunicodechar}
\newunicodechar{→}{\ensuremath{\rightarrow}}

\setmainfont{Helvetica Neue}
\setsansfont{Helvetica Neue}
\setlength{\parindent}{0pt}
\setlength{\parskip}{6pt}
\hyphenpenalty=200
\exhyphenpenalty=200
\pretolerance=400
\tolerance=2000
\setlength{\emergencystretch}{1em}
\renewcommand{\arraystretch}{1.25}
\setlist{leftmargin=5mm,itemsep=1mm,topsep=1mm}
\newcolumntype{Y}{>{\RaggedRight\arraybackslash}X}

\definecolor{milcMagenta}{HTML}{E6007E}
\definecolor{milcVino}{HTML}{5A0038}
\definecolor{milcTurquesa}{HTML}{0093A5}
\definecolor{milcVerde}{HTML}{58A923}
\definecolor{milcMostaza}{HTML}{E5B400}
\definecolor{milcOcre}{HTML}{B86600}
\definecolor{milcNegro}{HTML}{241816}
\definecolor{milcGris}{HTML}{F7F7F4}
\definecolor{milcLinea}{HTML}{E8E2DD}
\definecolor{milcGradePrimary}{HTML}{0066FF}
\definecolor{milcGradeDark}{HTML}{003D99}
\definecolor{milcGradeSoft}{HTML}{D6E8FF}
\definecolor{milcGradeAccent}{HTML}{CFE7FF}
\definecolor{milcGradeText}{HTML}{FFFFFF}
\color{milcNegro}

\pagestyle{fancy}
\fancyhf{}
\lhead{\small Tecnología e Informática · Grado 6}
\rhead{\small Guía 23 · MILC v3}
\cfoot{\small\thepage}
\renewcommand{\headrulewidth}{0pt}

\newtcolorbox{titlebox}[1]{
  enhanced,colback=#1,colframe=#1,arc=7mm,boxrule=0pt,
  left=7mm,right=7mm,top=3mm,bottom=3mm,
  before skip=8pt,after skip=8pt,
  fontupper=\sffamily\bfseries\Large\color{white}
}

\newtcolorbox{softbox}[3]{
  enhanced,colback=#2,colframe=#1,borderline west={4pt}{0pt}{#1},
  arc=2mm,boxrule=.4pt,left=4mm,right=4mm,top=3mm,bottom=3mm,
  before skip=6pt,after skip=8pt,title={#3},coltitle=white,
  colbacktitle=#1,fonttitle=\sffamily\bfseries,fontupper=\RaggedRight,
  attach boxed title to top left={xshift=3mm,yshift=-2mm},
  boxed title style={arc=2mm, boxrule=0pt}
}

\newtcolorbox{drawbox}[2]{
  enhanced,colback=white,colframe=#1,arc=4mm,boxrule=1pt,
  left=5mm,right=5mm,top=4mm,bottom=4mm,before skip=8pt,after skip=8pt,
  title={#2},coltitle=white,colbacktitle=#1,fonttitle=\sffamily\bfseries,
  fontupper=\RaggedRight,
  attach boxed title to top left={xshift=4mm,yshift=-2mm},
  boxed title style={arc=3mm, boxrule=0pt}
}

\newtcolorbox{infoband}[1]{
  enhanced,colback=#1!8,colframe=#1!50,arc=2mm,boxrule=.5pt,
  left=4mm,right=4mm,top=2mm,bottom=2mm,before skip=4pt,after skip=4pt,
  fontupper=\small\RaggedRight
}

% Puente narrativo entre secciones (contrato editorial Conéctate).
% Da continuidad: "Acabas de X. Ahora vas a Y porque Z."
% Visualmente: banda izquierda en color del grado, texto en itálica pequeña.
\newtcolorbox{puentebox}{
  enhanced,colback=milcGradeSoft,colframe=milcGradeDark,
  borderline west={3pt}{0pt}{milcGradeDark},
  arc=0mm,boxrule=0pt,left=5mm,right=4mm,top=2.5mm,bottom=2.5mm,
  before skip=4pt,after skip=8pt,
  fontupper=\itshape\small\color{milcNegro}\RaggedRight
}
% La flecha va en modo matematico a proposito: Helvetica Neue NO tiene el glifo
% U+2192, asi que \textrightarrow se compilaba a NADA («Missing character: There
% is no → in font Helvetica Neue») y el puente salia sin flecha. En modo
% matematico la flecha viene de la fuente matematica y si se dibuja.
\newcommand{\puente}[1]{%
  \begin{puentebox}%
    \textcolor{milcGradeDark}{$\rightarrow$}\hspace{2mm}#1%
  \end{puentebox}%
}

\newcommand{\linead}[1]{\noindent\color{milcLinea}\rule{#1}{0.5pt}\color{milcNegro}}
\newcommand{\respuesta}[1]{\par\vspace{2mm}\foreach \n in {1,...,#1}{\noindent\color{milcLinea}\rule{\linewidth}{0.5pt}\par\vspace{5mm}}\color{milcNegro}}
\newcommand{\checkbox}{\textcolor{milcGradeDark}{\fbox{\phantom{x}}}\hspace{5pt}}

\newcommand{\phasecard}[4]{
\begin{tcolorbox}[enhanced,colback=#3,colframe=#2,arc=3mm,boxrule=.5pt,height=3.05cm,valign=center,halign=center,left=1.5mm,right=1.5mm,top=2mm,bottom=2mm]
{\sffamily\bfseries\fontsize{22}{24}\selectfont\textcolor{#2}{#1}}\par
{\sffamily\bfseries\small #4}
\end{tcolorbox}
}

% ── Recursos visuales de la guía ──────────────────────────────────────
% Declarados en el bloque `recursos:` del YAML y emitidos por el builder.
% \guiaFigura   → imagen rasterizada o PDF (foto, carta celeste, montaje).
% \guiaDiagrama → fuente TikZ/circuitikz (.tex) incrustada como vector:
%                 es la vía para circuitos y esquemáticos editables.
% #1 = ruta relativa al .tex · #2 = pie (puede ir vacío)
\newcommand{\guiaFigura}[2]{%
\begin{center}
\includegraphics[width=0.88\linewidth,height=0.38\textheight,keepaspectratio]{#1}\\[1.5mm]
{\footnotesize\itshape\textcolor{milcNegro!75}{#2}}
\end{center}
}

\newcommand{\guiaDiagrama}[2]{%
\begin{center}
\input{#1}\\[1.5mm]
{\footnotesize\itshape\textcolor{milcNegro!75}{#2}}
\end{center}
}

\begin{document}
\thispagestyle{empty}
\begin{tikzpicture}[remember picture,overlay]
  \fill[white] (current page.south west) rectangle (current page.north east);
  \path[fill=milcGradePrimary,rounded corners=16mm]
    ([xshift=.55cm,yshift=-.55cm]current page.north west) rectangle
    ([xshift=-.55cm,yshift=3.65cm]current page.south east);

  \node[anchor=north west,text=milcGradeText] at ([xshift=2.0cm,yshift=-1.85cm]current page.north west)
    {\sffamily\bfseries\fontsize{14}{16}\selectfont GRADO};
  \node[anchor=north west,text=milcGradeText,opacity=.82] at ([xshift=2.0cm,yshift=-2.75cm]current page.north west)
    {\sffamily\bfseries\fontsize{9}{11}\selectfont SERIE GUÍAS · TECNOLOGÍA E INFORMÁTICA};

  \begin{scope}[shift={([xshift=-3.05cm,yshift=-2.55cm]current page.north east)}]
    \fill[white,opacity=.80] (-.62,-.52) rectangle (.62,.52);
    \draw[milcGradeText,line width=1pt] (-.62,-.52) rectangle (.62,.52);
    \fill[milcGradeDark] (-.42,-.38) rectangle (-.22,.06);
    \fill[milcGradeAccent] (-.10,-.38) rectangle (.10,.24);
    \fill[milcGradePrimary!60!black] (.22,-.38) rectangle (.42,.38);
  \end{scope}

  \node[anchor=west,text=milcGradeText] at ([xshift=1.9cm,yshift=-12.85cm]current page.north west)
    {\sffamily\bfseries\fontsize{124}{124}\selectfont 6};
  % El circulito de grado (°) solo aplica a las guias de grado («11°»); en
  % Bebras y Semillero el numero grande es un momento/modulo, no un grado.
  \draw[milcGradeText,line width=1.7pt]
    ([xshift=4.52cm,yshift=-11.68cm]current page.north west) circle (.13cm);

  \node[anchor=west,text=milcGradeText,text width=8.7cm,align=left] at ([xshift=10.35cm,yshift=-11.6cm]current page.north west)
    {
     {\sffamily\bfseries\fontsize{19}{22}\selectfont Formato\\básico ---\\darle\\voz visual\\a mi texto}\\[4mm]
     {\sffamily\bfseries\fontsize{23}{26}\selectfont Guía 23-6-TIC}\\[2mm]
     {\sffamily\fontsize{13}{16}\selectfont Período 3 · Procesador de texto e Internet}\\[5mm]
     {\sffamily\bfseries\fontsize{11}{13}\selectfont Institución Educativa Sor María Juliana}\\[1mm]
     {\sffamily\fontsize{10}{12}\selectfont Autor: Dr. Álvaro Cárdenas Orozco}
    };

  \path[fill=milcGradeSoft,rounded corners=7mm]
    ([xshift=1.35cm,yshift=.85cm]current page.south west) rectangle
    ([xshift=-1.35cm,yshift=4.25cm]current page.south east);
  \draw[milcGradeDark,line width=.65pt,rounded corners=7mm]
    ([xshift=1.35cm,yshift=.85cm]current page.south west) rectangle
    ([xshift=-1.35cm,yshift=4.25cm]current page.south east);
  \node[anchor=center,text=milcNegro,text width=17.0cm,align=center] at ([yshift=2.55cm]current page.south)
    {\sffamily\fontsize{10}{12}\selectfont Metodología MILC · Apertura ancestral · Pensamiento computacional · Triángulo Dussel-Estoicismo-Floridi\\
    \textcolor{milcGradeDark}{\bfseries Producto:} Tu documento de la S2 (los 3 párrafos) ahora con \textbf{formato profesional aplicado}: título en \emph{negrita centrada con tamaño 18}, párrafos en \emph{tamaño 12 justificados con sangría en primera línea}, \textbf{una palabra clave subrayada} en cada párrafo, \textbf{una idea importante en negrita} por párrafo, y \textbf{color sutil} en el título (azul oscuro). Más \textbf{una tabla en el cuaderno} con los 8 atajos de teclado que usaste, con dibujo de la combinación.};
\end{tikzpicture}
\mbox{}
\newpage

\section*{Apertura: Formato básico --- darle voz visual a lo que escribo}

\begin{softbox}{milcGradeDark}{milcGradeSoft}{Saber ancestral}
\textbf{Doña Mercedes la maestra rural de la vereda La Plata de Cartago decía: ``La letra entra con buena presentación.''} Cuando los niños campesinos le entregaban sus tareas, ella miraba 2 cosas: \textbf{el contenido} (¿qué dice?) y \textbf{la presentación} (¿se entiende a primera vista?). Una tarea con letra clara, márgenes respetados, títulos subrayados y palabras importantes resaltadas comunicaba al lector que el alumno se había tomado el trabajo en serio. Una tarea descuidada, sin importar qué tan inteligente fuera el contenido, daba pereza leer. \textbf{Doña Mercedes no era exigente por capricho}: enseñaba que \emph{la presentación es la primera forma de respeto al lector}. Hoy escribes en digital, pero \textbf{la regla sigue intacta}: dar formato no es decoración --- es respeto. Y los procesadores de texto te dan herramientas para hacerlo en segundos.
\end{softbox}

\begin{softbox}{milcTurquesa}{milcGris}{Saber contemporáneo}
El \textbf{formato} es la apariencia visual del texto: tipo de letra (\emph{fuente}), tamaño, color, alineación, énfasis. Cuando un documento está bien formateado, el lector \emph{ve la estructura antes de leer}: títulos arriba, lo importante resaltado, párrafos ordenados. \textbf{Los 5 elementos básicos de formato} son: \textbf{(1) Fuente}: tipo de letra (Arial, Calibri, Times New Roman). \textbf{(2) Tamaño}: 10-12 para texto normal, 14-18 para títulos. \textbf{(3) Estilo}: negrita (B), cursiva (I), subrayado (U). \textbf{(4) Alineación}: izquierda, centrada, derecha, justificada. \textbf{(5) Color}: el texto en negro siempre, los títulos pueden ir en azul oscuro o gris oscuro. Aprender estos 5 elementos te permite \textbf{transformar un texto plano en un documento profesional} en 2 minutos.
\end{softbox}

\begin{softbox}{milcMostaza}{milcGris}{Pregunta-puente}
\textit{Dos compañeros entregan una tarea con el mismo contenido. Uno la entrega \emph{con todo el texto en una sola masa, sin títulos, sin nada destacado}. El otro la entrega \emph{con título centrado en negrita, palabras importantes destacadas, párrafos separados}. ¿\textbf{Cuál crees que se va a leer primero}? ¿Por qué?}
\end{softbox}

\begin{softbox}{milcVerde}{milcGris}{Saber y saber hacer}
Al terminar la clase: (1) podrás \textbf{identificar} los 5 elementos básicos de formato; (2) sabrás \textbf{aplicar} los 8 atajos de teclado más útiles; (3) podrás \textbf{evaluar} si un documento tiene buen formato o no; (4) habrás \textbf{aplicado} formato profesional a tu documento de la S2.
\end{softbox}



\small
\begin{tabularx}{\textwidth}{>{\bfseries}p{3.0cm}Y}
\rowcolor{milcGradePrimary!12}Autor & Dr. Álvaro Cárdenas Orozco\\
DBA & Produce contenidos digitales y valida información de fuentes web (MEN, T\&I 6°)\\
Referentes & Saber ancestral de la maestra rural · Lectura crítica · Soberanía informacional\\
Producto final & Tu documento de la S2 (los 3 párrafos) ahora con \textbf{formato profesional aplicado}: título en \emph{negrita centrada con tamaño 18}, párrafos en \emph{tamaño 12 justificados con sangría en primera línea}, \textbf{una palabra clave subrayada} en cada párrafo, \textbf{una idea importante en negrita} por párrafo, y \textbf{color sutil} en el título (azul oscuro). Más \textbf{una tabla en el cuaderno} con los 8 atajos de teclado que usaste, con dibujo de la combinación.\\
\end{tabularx}
\normalsize

\puente{Hoy haces 4 cosas: identificas los 5 elementos de formato, memorizas los 8 atajos clave, aplicas formato a tu documento de la S2, y armas tu tabla de atajos.}

\begin{titlebox}{milcGradeDark}
Ruta MILC: así vamos a aprender
\end{titlebox}
\begin{tabularx}{\textwidth}{>{\centering\arraybackslash}X>{\centering\arraybackslash}X>{\centering\arraybackslash}X>{\centering\arraybackslash}X}
\phasecard{1}{milcVerde}{milcGris}{Escucha\\[-1mm]\small Observo, escucho y pregunto.} &
\phasecard{2}{milcTurquesa}{milcGris}{Sistematización\\[-1mm]\small Pensamiento computacional.} &
\phasecard{3}{milcMagenta}{milcGris}{Praxis\\[-1mm]\small Construyo y aplico.} &
\phasecard{4}{milcVino}{milcGris}{Evaluación\\[-1mm]\small Reflexión liberadora.}
\end{tabularx}


\puente{Empezamos comparando 2 versiones del mismo texto: una sin formato y una con formato.}

\begin{titlebox}{milcVerde}
Fase 1 · Escucha
\end{titlebox}

\begin{softbox}{milcVerde}{milcGris}{Escena para escuchar}
\textbf{Actividad 1 · ANALIZA --- Comparar 2 versiones del mismo texto} (10 min · individual). Lee estos 2 textos. Tienen exactamente el mismo contenido, solo cambia el formato.\\[2mm] \textbf{VERSIÓN A (sin formato):} \emph{``Cartago es una ciudad del valle del cauca colombia esta en el norte del departamento se fundó en 1540 por jorge robledo es famosa por su cultura cafetera tiene un parque central llamado parque de bolívar la gente es amable y cálida vivir allí es agradable''}.\\[2mm] \textbf{VERSIÓN B (con formato):} título \emph{``Cartago, Valle''} grande y centrado; un párrafo con palabras clave en negrita (\emph{Cartago}, \emph{1540}, \emph{Jorge Robledo}), una palabra subrayada (\emph{cafetera}), todo el texto justificado con sangría. \textbf{En el cuaderno responde:} \emph{(a) ¿Cuál se entiende más rápido?} \emph{(b) ¿Cuál parece más profesional?} \emph{(c) ¿Cuál vas a leer primero si tienes prisa?}.
\end{softbox}

\checkbox Leí las 2 versiones con calma.\\
\checkbox Respondí las 3 preguntas con honestidad.\\
\checkbox Note la diferencia clara entre las dos.

\begin{infoband}{milcVerde}


\textbf{Lo que va al cuaderno (Actividad 1 · ANALIZA).} Título: «Actividad 1 --- Texto con y sin formato». \textbf{Formato:} 3 respuestas cortas a las 3 preguntas. \textbf{Extensión:} 3 líneas. \textbf{Sabes que terminaste cuando} entiendes intuitivamente por qué el formato importa.
\end{infoband}

\puente{Después aprendes los 5 elementos + los 8 atajos de teclado más útiles.}

\begin{titlebox}{milcTurquesa}
Fase 2 · Sistematización con pensamiento computacional
\end{titlebox}

\begin{softbox}{milcTurquesa}{milcGris}{Cuatro pilares aplicados al tema}
Lo que viste es real: \textbf{el mismo contenido se siente totalmente distinto según el formato}. Ahora aprenderás los \textbf{5 elementos} para aplicar y los \textbf{8 atajos de teclado} que te harán rápido.
\end{softbox}

\small
\begin{tabularx}{\textwidth}{>{\bfseries\RaggedRight\arraybackslash}p{3.6cm}Y}
\rowcolor{milcTurquesa!90}\color{white}Pilar & \color{white}\bfseries Aplicación al tema\\
1. Descomposición & \textbf{Elemento 1 --- Fuente (tipo de letra).} La que viene por defecto en Word es \emph{Calibri}. Otras populares: \emph{Arial} (sencilla), \emph{Times New Roman} (clásica para informes), \emph{Verdana} (legible en pantalla). \textbf{Regla:} usa \emph{máximo 2 fuentes} en un documento --- una para títulos, otra para texto. Más de 2 = caos visual. \textbf{Elemento 2 --- Tamaño.} \emph{12 puntos} es el estándar para texto. \emph{10-11} para notas pequeñas. \emph{14-18} para títulos. \emph{20+} solo para portadas. \emph{Regla:} no cambies el tamaño dentro del mismo párrafo.\\
2. Reconocimiento de patrones & \textbf{Elemento 3 --- Estilo (negrita, cursiva, subrayado).} \textbf{Negrita (Ctrl+N o Ctrl+B)}: para \emph{ideas importantes}, no para párrafos enteros. \emph{Cursiva (Ctrl+K o Ctrl+I)}: para citas, palabras en otro idioma, títulos de libros. \emph{Subrayado (Ctrl+S o Ctrl+U)}: para palabras clave o términos técnicos. \textbf{Regla:} usa cada estilo en máximo \emph{2-3 palabras por párrafo}. Más, pierde el énfasis.\\
3. Abstracción & \textbf{Elemento 4 --- Alineación.} \emph{Izquierda (Ctrl+Q)}: la usual para texto, ideal para correos y cartas. \emph{Centrada (Ctrl+E)}: solo para títulos y citas destacadas. \emph{Derecha (Ctrl+D)}: para fechas y firmas. \emph{Justificada (Ctrl+J)}: el texto se estira para que llegue a los 2 márgenes, da aspecto profesional. \textbf{Regla:} usa \emph{justificada} para documentos formales, \emph{izquierda} para informales.\\
4. Algoritmo conceptual & \textbf{Elemento 5 --- Color.} \emph{Negro}: SIEMPRE para el texto del cuerpo. Cambiar el color del texto cuerpo a otro color hace ver el documento infantil. \emph{Azul oscuro, gris oscuro, vino tinto:} aceptables solo para títulos. \textbf{Nunca:} amarillo, rosado fucsia, verde fluorescente en texto principal. \textbf{Los 8 atajos clave para memorizar:} (1) \texttt{Ctrl+N} = negrita; (2) \texttt{Ctrl+K} = cursiva; (3) \texttt{Ctrl+S} = subrayado; (4) \texttt{Ctrl+E} = centrar; (5) \texttt{Ctrl+J} = justificar; (6) \texttt{Ctrl+G} = guardar; (7) \texttt{Ctrl+Z} = deshacer; (8) \texttt{Ctrl+C / Ctrl+V} = copiar y pegar.\\
\end{tabularx}
\normalsize

\begin{drawbox}{milcTurquesa}{5 elementos + 8 atajos}
\textbf{5 elementos:} fuente · tamaño · estilo · alineación · color. \\[1mm]
\textbf{8 atajos:} \\[1mm]
Ctrl+N (negrita) · Ctrl+K (cursiva) · Ctrl+S (subrayado) \\[1mm]
Ctrl+E (centrar) · Ctrl+J (justificar) · Ctrl+G (guardar) \\[1mm]
Ctrl+Z (deshacer) · Ctrl+C/V (copiar/pegar).
\end{drawbox}

\begin{softbox}{milcMostaza}{milcGris}{Errores comunes y su corrección}
\textbf{Errores de formato que dejan mal un documento.} (1) \emph{Todo en mayúsculas}: parece que gritas y se lee peor. (2) \emph{Muchas fuentes mezcladas}: parece collage de niño. (3) \emph{Colores chillones}: rosado, amarillo fluorescente, verde fluorescente. (4) \emph{Negrita en todo el texto}: pierde el énfasis (si todo es importante, nada es importante). (5) \emph{Tamaños distintos entre párrafos}: parece error de copy-paste. (6) \emph{Cursivas para enfatizar lo importante}: la cursiva es para citas y palabras extranjeras, NO para énfasis fuerte (esa es la negrita).
\end{softbox}

\begin{infoband}{milcTurquesa}


\textbf{Lo que va al cuaderno (Actividad 2 · EXPLICA).} Título: «Actividad 2 --- 5 elementos + 8 atajos». \textbf{Formato:} 2 bloques. Bloque A: tabla de 5 filas con \emph{elemento + regla principal}. Bloque B: tabla de 8 filas con \emph{atajo + qué hace + dibujo de la combinación de teclas} (un cuadrito por tecla). \textbf{Extensión:} 5 + 8 filas. \textbf{Sabes que terminaste cuando} podrías formatear un texto sin necesidad de mirar la barra de menú.
\end{infoband}

\puente{Con los atajos en mente, abres tu documento de la S2 y le das formato profesional.}

\begin{titlebox}{milcMagenta}
Fase 3 · Praxis con pensamiento computacional
\end{titlebox}

\begin{softbox}{milcMagenta}{milcGris}{Construyendo el producto con los cuatro pilares}
\textbf{Actividad 3 · APLICA --- Formatear tu documento de la S2} (28 min · individual). Vas a abrir el documento de \emph{Mi primer texto digital} que creaste en S2 y le aplicarás formato profesional usando los 8 atajos.
\end{softbox}

\small
\begin{tabularx}{\textwidth}{>{\bfseries\RaggedRight\arraybackslash}p{3.6cm}Y}
\rowcolor{milcMagenta!90}\color{white}Pilar & \color{white}\bfseries Aplicación al producto\\
Descomposición & \textbf{Paso 1 --- Abre tu documento.} Ve a la carpeta donde lo guardaste y haz doble clic. Si no lo encuentras, créalo de nuevo siguiendo S2 con 3 párrafos sobre un tema tuyo.\\
Patrones & \textbf{Paso 2 --- Formatea el título.} Selecciona el título (mantener clic izquierdo desde el inicio del título y arrastrar hasta el final). Aplica: \textbf{Ctrl+N (negrita)}, \textbf{Ctrl+E (centrar)}, cambia el \textbf{tamaño a 18} (en el cuadro de tamaño de la barra de arriba), y cambia el \textbf{color a azul oscuro} (en el menú de color de texto, escoge un azul que no sea fluorescente).\\
Abstracción & \textbf{Paso 3 --- Formatea los párrafos.} Selecciona los 3 párrafos completos. Aplica: \textbf{Ctrl+J (justificar)}, asegúrate de que estén en \textbf{tamaño 12}, fuente \emph{Calibri} o \emph{Arial}, color \emph{negro}.\\
Algoritmo de armado & \textbf{Paso 4 --- Resalta lo importante.} En cada uno de los 3 párrafos: identifica \textbf{1 palabra clave} y subráyala (selecciónala + Ctrl+S). Identifica \textbf{1 idea importante} (3-5 palabras) y ponla en negrita (selecciónala + Ctrl+N). En total tu documento tendrá 3 subrayados y 3 negritas. Ni más ni menos --- el énfasis se pierde si abusas.\\
\end{tabularx}
\normalsize

\begin{drawbox}{milcMagenta}{Antes de mostrar tu documento formateado}
\checkbox El título está en negrita, centrado, tamaño 18, color azul oscuro. \\[1mm]
\checkbox Los 3 párrafos están justificados, tamaño 12, color negro. \\[1mm]
\checkbox Hay 1 palabra subrayada en cada párrafo (3 en total). \\[1mm]
\checkbox Hay 1 idea en negrita en cada párrafo (3 en total). \\[1mm]
\checkbox La primera línea de cada párrafo tiene sangría. \\[1mm]
\checkbox El documento está guardado con el mismo nombre.
\end{drawbox}

\begin{softbox}{milcMagenta}{milcGris}{Plantilla de trabajo}
\textbf{Estructura del documento formateado.} \textit{Título:} negrita + centrado + tamaño 18 + azul oscuro. \textit{Cada párrafo:} justificado + tamaño 12 + negro + sangría primera línea. \textit{En cada párrafo:} 1 palabra subrayada + 1 idea en negrita. \textit{Atajos usados:} Ctrl+N, Ctrl+S, Ctrl+E, Ctrl+J, Ctrl+G.
\end{softbox}

\begin{infoband}{milcMagenta}


\textbf{Lo que va al cuaderno (Actividad 3 · APLICA).} Título: «Actividad 3 --- Mi documento con formato profesional». \textbf{Formato:} captura impresa del documento formateado (si tienes acceso a impresora) o boceto en el cuaderno + tabla de los 8 atajos con dibujo de la combinación de teclas + reflexión de 2 líneas comparando antes/después. \textbf{Extensión:} 1 página entera del cuaderno. \textbf{Sabes que terminaste cuando} si miras tu documento de la S2 al lado del de hoy, sientes el salto de calidad.
\end{infoband}

\puente{El producto es ese documento formateado + la tabla de atajos en el cuaderno.}

\begin{titlebox}{milcMostaza}
Producto de la sesión
\end{titlebox}
\textbf{Documento formateado profesionalmente + tabla de 8 atajos · firmado}.
\begin{softbox}{milcMostaza}{milcGris}{Criterios de calidad}
(1) El título tiene negrita, centrado, tamaño 18 y color azul oscuro. (2) Los párrafos están justificados con sangría en primera línea. (3) Hay énfasis controlado (1 subrayado + 1 negrita por párrafo). (4) El estudiante usó los atajos de teclado (no solo el menú de la barra). (5) La tabla de atajos está en el cuaderno con dibujo de la combinación.
\end{softbox}

\puente{Antes de salir, comparas tu documento antes y después: ¿se siente más profesional ahora?}

\begin{titlebox}{milcVino}
Fase 4 · Evaluación liberadora — cinco dimensiones
\end{titlebox}

\begin{softbox}{milcVino}{milcGris}{Sentido de la evaluación}
Evaluar no es solo revisar una nota. Es reconocer qué aprendí, qué cambió en mí, qué emociones manejé, cómo me relaciono con la ciudad y la comunidad, y qué le dejaré a quien viene detrás.
\end{softbox}

\textbf{Mi reflexión liberadora en cinco dimensiones.} Completa cada frase iniciada:

\small
\begin{tabularx}{\textwidth}{>{\bfseries\RaggedRight\arraybackslash}p{3.4cm}Y>{\RaggedRight\arraybackslash}p{4.6cm}}
\rowcolor{milcVino}\color{white}Dimensión & \color{white}\bfseries Mi respuesta (frame starter) & \color{white}\bfseries Algo que descubrí\\
1. Personal & Lo que más cambió en mí fue \linead{6.5cm} & \linead{4.2cm}\\
2. Emocional & La emoción más fuerte fue \linead{2.5cm} y la regulé \linead{3cm} & \linead{4.2cm}\\
3. Ciudadana & En democracia esto importa porque \linead{6cm} & \linead{4.2cm}\\
4. Local (Cartago) & En mi barrio sería útil para \linead{6.5cm} & \linead{4.2cm}\\
5. Intergeneracional & A mi abuela/abuelo le diría \linead{6.5cm} & \linead{4.2cm}\\
\end{tabularx}
\normalsize

\begin{infoband}{milcVino}
\textbf{Para la próxima sustentación.} Vuelve a esta tabla en seis meses. Verás que las respuestas cambian — no porque lo aprendido cambie, sino porque tú cambias. La evaluación liberadora es una conversación contigo a través del tiempo.
\end{infoband}


\puente{Cerramos con 3 voces que te ayudan a entender por qué dar formato es un acto de respeto al lector.}

\begin{titlebox}{milcGradeDark}
Triángulo de pensamiento · cierre filosófico
\end{titlebox}

\begin{softbox}{milcVino}{milcGris}{Dussel · lente del nosotros}
\textit{````El formato es el primer abrazo que le das a quien te va a leer. Una página bien presentada dice: te respeto.''''}

Cuando dejas tu documento sin formato, le dices al lector: \emph{``arréglatelas para entenderme''}. Cuando lo formateás bien, le dices: \emph{``te facilito el camino''}. Doña Mercedes, en la escuela rural, exigía buena presentación porque sabía que \textbf{leer es esfuerzo}: si tú no ayudas con formato, el lector tiene que poner el doble. Por eso \textbf{el formato es ética}, no decoración.

\textbf{Cuando entrego un documento, ¿le facilito el camino al lector o lo dejo solo?}
\end{softbox}
\linead{\linewidth}\\[1mm]
\linead{\linewidth}

\begin{softbox}{milcMostaza}{milcGris}{Estoicismo · Marco Aurelio · cuidado interior}
\textit{````Lo que está ordenado por fuera ayuda a la mente a ordenarse por dentro.''''}

Cuando \textbf{formateas bien un documento}, no solo ayudas al lector: también te ayudas a ti. Tu cerebro \emph{ve} la estructura y \emph{piensa} más claro. Si revisas algo desordenado, tu mente se confunde; si revisas algo ordenado, encuentras errores y mejoras rápido. El formato es \emph{disciplina externa que produce claridad interna}.

\textbf{¿Cómo afecta a mi propio pensamiento el desorden visual de un documento?}
\end{softbox}
\linead{\linewidth}\\[1mm]
\linead{\linewidth}

\begin{softbox}{milcTurquesa}{milcGris}{Floridi · lente de la infoesfera}
\textit{````Un documento digital bien formateado es un acto cívico: respeta el tiempo de quien lo leerá, sea uno o mil.''''}

Lo que escribes en digital puede ser leído por \textbf{muchas personas}: tu profe, tus papás, tus compañeros, futuros colegas, lectores que no conoces. Si tu documento está bien formateado, \textbf{respetas el tiempo de cada uno de ellos}. Si no, los obligas a perder tiempo descifrando. \emph{Formatear bien es ciudadanía digital}: cuidar el tiempo común.

\textbf{¿Mis documentos respetan el tiempo de quien va a leerlos?}
\end{softbox}
\linead{\linewidth}\\[1mm]
\linead{\linewidth}

\begin{titlebox}{milcVino}
Compromiso final
\end{titlebox}

\small
\begin{tabularx}{\textwidth}{>{\RaggedRight\arraybackslash}p{3.0cm}YYY}
\rowcolor{milcVino}\color{white}\bfseries Criterio & \color{white}\bfseries Lo logré & \color{white}\bfseries En proceso & \color{white}\bfseries Necesito apoyo\\
Comprensión del tema & Explico ideas centrales con mis palabras. & Reconozco ideas pero falta precisión. & Necesito repasar conceptos base.\\
Pensamiento computacional & Apliqué los 4 pilares al diseñar mi producto. & Apliqué algunos pilares. & Necesito releer la fase 2.\\
Aplicación y producto & Mi producto cumple los criterios y se puede revisar. & Producto funcional con ajustes pendientes. & Aún en construcción.\\
Reflexión 5 dimensiones & Respondí cada dimensión con honestidad. & Respondí algunas con detalle. & Mis respuestas son muy generales.\\
Triángulo de pensamiento & Conecté las 3 voces con mi trabajo real. & Conecté al menos una. & No relaciono las voces con mi proceso.\\
\end{tabularx}
\normalsize

\textbf{Hoy entendí que:}\\[1mm]
\linead{\linewidth}\\[1mm]
\linead{\linewidth}

\textbf{Mi compromiso para la próxima sesión es:}\\[1mm]
\linead{\linewidth}\\[1mm]
\linead{\linewidth}

\begin{softbox}{milcMostaza}{milcGris}{Cita de despedida · Marco Aurelio}
\textit{``El obstáculo en el camino se vuelve el camino. Lo que estorba al camino se vuelve camino.''}\\[1mm]
Cada error de hoy, bien observado, se convierte en pieza del camino que pisarás mañana.
\end{softbox}

\small
\begin{tabularx}{\textwidth}{>{\bfseries}p{3.6cm}Y}
\rowcolor{milcGradeDark!12}Nombre completo: & \linead{10cm}\\
Fecha: & \linead{4cm} \quad Curso/grupo: \linead{4cm}\\
Mi firma: & \linead{9cm}\\
\end{tabularx}
\normalsize

\begin{infoband}{milcGradeDark}
\textbf{Para la próxima sesión.} Esta semana, formateás cualquier tarea que entregues en digital usando los 8 atajos. No es opcional: es tu nueva firma. La próxima clase aprendes \textbf{listas, viñetas y numeración} --- las herramientas para mostrar información estructurada. Vas a poder transformar un párrafo confuso en una lista limpia que se entiende en 5 segundos.
\end{infoband}

\end{document}
